\documentclass[2pt]{book}
\setcounter{secnumdepth}{3}
\setcounter{tocdepth}{3}
\usepackage{graphicx}
\usepackage{html}
\usepackage{hyperref}
\usepackage{lmodern}
\usepackage{tocloft}
\usepackage{svg}
\usepackage{listings}
\setlength{\cftpartindent}{0em}
\setlength{\cftchapindent}{0em}
\setlength{\cftsecindent}{0em}
\setlength{\cftsubsecindent}{0em}
\setlength{\cftsubsubsecindent}{0em}
\graphicspath{{Images/}}
\author{Doug Royer / RiverExplorer™}
%\date{13-MAR-2025}
\title{Phoenix: Server and Client Implementation Design}
%
% Define Abstract ..
%
\newenvironment{abstract}
{%
  \small
  \begin{center}
    \bfseries{abstract}
  \end{center}
}
{}
\lstset{
  language=C++,
  basicstyle=\ttfamily,
  keywordstyle=\color{blue},
  commentstyle=\color{brown},
  stringstyle=\color{red},
  numbers=left,
  numberstyle=\tiny,
  stepnumber=1,
  numbersep=5pt,
  backgroundcolor=\color{gray!10!white},
  showspaces=false,
  showstringspaces=false,
  showtabs=false,
  frame=single,
  tabsize=2,
  captionpos=b,
  breaklines=true,
  breakatwhitespace=false,
  title=\lstname,
  escapeinside={\%*}{*)},
  morekeywords={*,...}
}

\begin{document}
\maketitle
\frontmatter
\begin{figure}
  \centering
  \includegraphics{RiverExplorerLLC-Logo-BLACK-131x50-Transparent.ps.png}
\end{figure}

\begin{abstract}
  \includegraphics{UnderConstruction.png}
  \begin{center}
    \Huge{\textcolor{red}{\textbf{THIS IS WORK IN PROGRESS}}}
  \end{center}
  \begin{center}
    \underline{\textit{ABSTRACT}}
  \end{center}
  Email and MIME messages account for one the largest volumes of data
  on the internet. The transfer of these MIME message has not had a
  major updated in decades.
  Part of the reason is that it is very important data and altering it
  takes a great deal of care and planning.

  This application transport can also transfer non-MIME data.
  It can be used as an XDR transport, or for opaque data (blobs
  of known or unknown data) transport.

  Another major concern is security and authentication.
  This proposal allows for existing authentication to continue to work.

  This is a MIME message transport that can facilitate the transfer
  of any kind of MIME message.
  Including email, calendaring, and text, image, or multimedia MIME
  messages. It can transfer multipart and simple MIME messages.

  The POP and IMAP protocols are overly chatty and now that the
  Internet can handle 8-bit transfers, there is no need for the
  overly complex text handling of messages.

  This proposal includes a sample implementation.
  (Github - Phoenix) Which also includes a gateway from this proposal
  to existing systems.
  Thunderbird and Outlook plugins are part of the sample implementation.
  A Linux, Windows DLL and .NET, and Android client library are
  part of the sample implementation

  Phoenix is a MIME Transport Protocol (MTP).
  Named Phoenix because it is a risen version version of the
  Lemonade project if the late 1990's.

  Unless otherwise specified, all of this code and text is original
  and writtern by the author.
\end{abstract}


\tableofcontents

\mainmatter
\chapter{Introduction}
\section{Introduction}
The client and the server send data to the network, and get
data from the network.
This \mbox{input/output} model (IO Model) works with
both the Phoenix client and Phoenix server.

\textbf{NOTE: It is possible that specific implementations allow
  each endpoint to act as a client and server.
  For the purposes of this specification, the client is
  the one that is originating actions and expects
  results or replies from the other endpoint}

This model support both threaded and non threaded
implementations.
For non-threaded systems the data is processed serially
or however the implementation handles both incoming
and outgoing data.
The data flow is the same for both threaded and non threaded
implementations.

\begin{figure}
  \centering
  \caption{I/O Model}
  \label{fig:IOModel}
  \includesvg{Generated/IO-00.svg}
\end{figure}

\begin{description}
\item[(A)Main Thread]\label{item:MainThread}
	The main thread processes data and formats it into
  IOVectorPacket objects that contains the command
  and references to the data that will be sent.
  Then stores the IOVectorPackt (figure: \ref{fig:IovPacket})
  into queue \textit{OQ}.
  The client \textit{Main Thread} is described in section \ref{sec:ClientMainThread}
  And the server \textit{Main Thread} is described in section \ref{sec:ServerMainThread}

\item[(B) Output Thread]
  This thread waits for the output queue \textit{OQ} to have
  content.
  It checks to make sure that the network is writable and not write blocked
  or otherwise busy.
  When it can proceed, it takes an entry out of the queue,
  and (C) XDR encodes it.
  The XDR stream is opened to the TLS encryption.
  The \textit{Output Thread} is described in section \ref{sec:OutputThread}
  
\item[(C) XDR Encoding]\label{item:XdrEncodingThread}
	This thread, takes one packet at at time
	out of the input queue,	converts it from the native computer
  to network compatible XDR format.
  It reads the IOVectorPacket data to an XDR stream and that streams is sent directly
  to the \textit{(D) TLS Encryption} code.
  \textit{XDR Encoding} is described in section \ref{sec:XdrEncode}

\item[(D) TLS Encryption]\label{item:TlsEncryption}
  This code takes the XDR encoded data from its input stream
  and encrypts it to the network.
  Which encrypts it according to the certificates supplied by the configuration.
  and sends it out to the network.
  \textit{TLS Encryption} is described in section \ref{sec:TlsIOOut}

\item[(E) Input Thread]\label{item:InputThread}
  This thread waits for any data coming from the network
  and feeds it into the \textit{TLS Decryption} code.
    
\item[(F) TLS Decryption]\label{item:TlsDecryptionThread}
  This code decrypts it on the fly, and streams it to \textit{(G) XDR Decoding} code.
  \textit{TLS Decryption} is described in section \ref{sec:TlsIOIn}

\item[(G) XDR DECODING]\label{item:XdrDecoding}
  This takes data from the \textit{TLS Decryption} code and decodes
  it back to native computer data format.
  \textit{XDR Decoding} is described in section \ref{sec:XdrDecode}
  
\item[(H) Dispatch Thread]\label{DispatchThread}
	This is the dispatch thread.
  It takes the decoded packet and dispatches it to the code that can handle the data.
  The \textit{Dispatch Thread} is described in section \ref{sec:DispatchThread}

\item[(OQ) Output Queue]
  This is a C++ \textit{std::deque} of I/O Vector Packets
  shown in figure \ref{fig:IovPacket}.
  
\item[(IQ) Input Queue]
  This is a C++ \textit{std::deque} of I/O Vector Packets
  shown in figure 1\ref{fig:IovPacket}.
  
\end{description}


This model is the core of both the client and server.

\chapter{Commands}
\section{Commands Overview}
Commands are all \ref{code:PacketBody}{ PacketBody} objects.

The command descriptions are broken down into related sections in order
to simplify the documentation.

The client connects to the server, and then the
server waits for a AUTHENTICATION command or a PRE
authentication CAPABILITY command from the client.

There are PRE and POST CAPABILITY commands.
The PRE CAPABILITY command have to do with establishing
an authenticated user or client connection.
The POST authentication CAPABILITY command informs
the other endpoint what it can do for the other endpoint
on behalf of this now authenticated user or client.

One or more commands are sent in each PacketBody
sent between the client and server.

There is no guarantee that any two clients and server
can authenticate or process any data.
The Phoenix protocol is an application level MIME transport
protocol and the two endpoints might not implement the
same MIME types.
For example one might do only email, and the other might do
only RSS feeds.
Other servers and clients may be capable of multiple tasks.

Commands can be documented commands or vendor commands.
A vendor may publicly document its extensions, or they
may elect to keep them proprietary and unavailable to other vendors.
Vendor commands are allowed because the purpose of this protocol
is to facilitate the transfer of MIME objects.
It is not to guarantee every implementation can transfer
any data to or through any other implementation.

\subsection{Authentication Commands}
AUTHENTICATION commands is the Authentication, authorization,
accounting (AAA) process between a client and server.

More than one authentication mechanisms may be required
by a server.
For example a server may require AUTHCERT\_TLS and AUTHCERT\_USER
or perhaps AUTHCERT\_TLS and AUTHMD5.

\subsubsection{New Client and Server Relationship}
When a client has not had a successful authentication to a
specific server, it sends its first packet as a PRE authentication
CAPABILITY packet informing the server of the authentication mechanisms
it supports.
The list is prioritized.
With the most preferred listed first.
And the least preferred ones later in the list.

Once the server gets this list, it picks the highest
priority one it supports then then sends the client
the servers PRE authentication CAPABILITY command listing only
that authentication mechanism.
The server may include any other PRE authentication CAPABILITY
items to the client it may have in that command.

If no common authentication mechanism are satisfactory,
the server sends the client a BYE
command and closes the connection.

If the client gets a BYE command, it then knows that it can
not communicate with this server and also closes its connection.

If the client gets the servers PRE authentication CAPABILITY
command, it attempts the agreed on authentication process
using any possible additional CAPABILITY description
restrictions or conditions provided in the server
PRE authentication CAPABILITY contents.

When the agreed on authentication mechanism works,
the server sends its POST authentication CAPABILITY
list to the client.

When the agreed on authentication mechanism fails,
the server a PRE authentication CAPABILITY command
to the client.
It can be the same as sent before, or modified.

A server may at any time block or temporally block an IP address
from attempting to start authentication again or too
quickly after a failure.

\subsubsection{Existing Client and Server Relationship}
When a client has had a relationship to a specific server
using a specific authentication mechanism it may bypass
the PRE authentication CAPABILITY communications and send an
AUTHENTICATION command when first connecting in order to reduce
the round trips needed to start processing other commands.
This is done by sending the AUTHENTICATION mechanism that had
previously worked with this server as the first command
in a PacketBody that also includes the clients PRE authentication
CAPABILITY command.
The first command is the AUTHENTICATION and the second included
command is the clients PRE authentication CAPABILITY command.
When this authentication succeeds, the server sends its POST
authentication CAPABILITY back to the client as the indicator
that it worked and the client is authenticated.

If this initial authentication fails, client is not yet authenticated.
And the server having already received the clients PRE authentication
CAPABILITY command, picks the highest priority authentication
mechanism they share and sends the client the servers PRE authentication
command containing any other CAPABILITY items it may need to share.

Then the authentication continues until success or failure.

When a server, after examining the received client PRE authentication
CAPABILITY command determines that they do not share any
common authentication, the server sends the client a BYE
command and closes the connection.
The BYE command at this step informs the client that they
share no common authentication mechanisms and also
closes the connection.

\subsubsection{AUTHANONYMOUS}
An anonymous authentication would be allowed on
servers that do not require a user to login to the system.
An example could be a \textit{Really Simple Syndication} (RSS)
feed that provides the latest product release notes available to anyone.

Server administration need to ensure that absolutely anyone
with network access has permission to connect to an AUTHANONYMOUS
server.

\subsubsection{AUTHCERT\_TLS}
TLS authentication is done when the client connects to the
server.

They each validate the TLS certificate and certificate chain
of the other.

If they validate, the client is authenticated to the server.

If they do not validate, the client may choose another authentication
mechanism or disconnect from the server.

A server may disconnect from the client if AUTHCERT\_TLS is required.

\subsubsection{AUTHCERT\_USER}
AUTHCERT\_USER is a secondary certificate that is used
after the initial TLS connection.

For example, if the user is travailing, they may have
a portable certificate tool or software that is used
to ensure that unknown or unfriendly hardware can not
view the communications.

The steps are:
\begin{enumerate}
\item Establish the TLS connection.
\item Using the user certificate, encrypt all communication
  from this point forward.
  Doubling the encryption process.
\end{enumerate}

A client that implements AUTHCERT\_USER would need a way
for a user to connect a device to the client in order
to allow the data flow to go through the user supplied hardware.
An example could be a specialty USB device that encrypts the data
between the network and the client.

\subsubsection{AUTHMD5}
MD5 authentication is used in many protocols.
The client, over a TLS connection sends the account name
and the MD5 hashed password to the server.

\subsubsection{BYE}
When the client or server receives a BYE command,
it has been informed that the other endpoint has closed
or is closing the connection and that no further communications
can take place on this connection.

\subsection{Capability Commands}
Capability commands are not a negotiation.
Each endpoint informs the other endpoint of the type
of transmissions it can process.

There is no guarantee that any two clients and server
can process any data.
The Phoenix protocol is an application level MIME transport
protocol and the two endpoints might not implement the
same MIME types.
For example one might do only email, and the other might do
only calendar feeds.
Other servers and clients may be capable of multiple tasks.

\subsubsection{CAPABILITY\_PRE}
A PRE authentication CAPABILITY is used before
the endpoints are authenticated and contain minimal
information.

At a minimum a PRE authentication capability includes
a prioritized list of authentication mechanisms it supports.
It also includes a VENDOR\_ID which is used to let each
endpoint know if they can transfer or perform vendor specific
data or communications.

\subsubsection{CAPABILITY\_POST}
A POST authentication CAPABILITY is sent after a successful
authentication informing the other endpoint of the functions
that it can perform.

A server may include the VENDOR\_ID value again, if after
authentication the server determines the client an perform
actions that are unique between the client and server.

A vendor may publicly document its extensions, or they
may elect to keep them proprietary and unavailable to other vendors.
\subsubsection{VENDOR\_ID}

\subsection{Control Commands}
These are housekeeping command for the Phoenix protocol.

\subsubsection{NOT\_SUPPORTED}
A NOT\_SUPPORTED command is sent as the reply to a command
not supported by the receiving endpoint.
The sequence number of the received command is sent in
the NOT\_SUPPORTED command to indicate which command
is not supported.

\subsubsection{PING}

The ping command is only sent when the client implementation has
determined it has waited too long for a command reply.
The ping command is only initiated from the client.
It is not valid for the server to send a ping command to a client.

A ping command must not be sent before a successful authentication.

The ping command MUST NOT be the first command sent to the server.
It should only be sent when the client implementation determines
it has waited too long for a reply.

If the server supports the ping command, then a PING capability is
sent in the CAPABILITY\_POST command.

Sometimes servers are unavailable and can go down.
A server could be down for maintenance, or in a shutdown mode.
It might limit the number of simultaneous connections.
It might be very busy.
The packets might not be making it to the server because of network
issues.

When a ping command is received by the server:
\begin{enumerate}
\item When the server did not send PING capability in the post
authentication capability list to the client.
The server ignores the PING command.
\item When the connection is not authenticated,
  The server ignores the PING command.
\item When the client is authenticated, and when the server is
  available for processing commands.
  Then the server replies with a ping reply with the same sequence number.
\end{enumerate}

If the server is alive and not available, the server will reply with
a NotSupported command, with its sequence number set to the sequence
number in the ping command.

If a connected and authenticated client has been waiting for a reply
or for some other reason needs to determine if the server is still
available.
It can send a ping command. If the server is still available, it
sends a ping reply.
If it is no longer available for any reason, it sends
a NotSupported reply.

A client MUST NOT send a ping command if it is waiting the results
of a previously sent ping command.
If the server is ignoring PING commands, a reply will never happen.

A client MUST NOT send a ping command more frequently than 90\% of
the SERVER\_TIMEOUT value that the server sent in the
CAPABILITY\_POST command.

Servers must give priority to ping commands.
If possible, reply as soon as it receives the command.

With servers that support PING, clients MUST NOT send any other
command while waiting for the PING reply.

The server MAY consider too many ping commands as a
malfunctioning or malicious client and terminate the connection.

Servers that are not threaded or can not reply to simultaneous
or overlapping commands, MUST NOT include PING in their post
authentication capability command.

\subsubsection{RESERVED\_CMD}
The RESERVED\_CMD is a placeholder for future extensions.
Its value is 0xffffffff.

\subsection{File and Folder Commands}
A folder is not necessarily a directory on the server.
It could be a logical folder in another object.
It might be a specific email message in a email folder container.
It could be a specific file in a zip or tar file.
It could be an entry in a database.

Within the Phoenix protocol a folder is simply a container
of any kind that can perform one or more of the folder
and file commands Phoenix supports.
Not all folders may support all commands.

Examples: Some may be read only.
Others might not allow them to be renamed.

\subsubsection{FILE\_CREATE}
Create a new file.
What is creates is a function of the task the
server is being asked to do:

Examples of what it could be doing:
\begin{enumerate}
\item When it is processing email messages, it could be creating
  a new email message.
  When sent to an outbound email server, it sends email.
  When sent to a user email server, it could be saving a draft.
\item When it is processing calendar messages, it could be creating
  a new calendar event, todo, or other calendar object.
\item It could be creating a new file on a disk.
\end{enumerate}

See the specific task specifications.

A FILE\_CREATE command includes the ID of the folder
to create.
And if applicable, a path to the new file.
The path separation character is the slash (/) 0x2f UTF-8 character.
\subsubsection{FILE\_COPY}

\subsubsection{FILE\_DELETE}
\subsubsection{FILE\_RENAME}
\subsubsection{FILE\_METADATA}
\subsubsection{FILE\_MOVE}
\subsubsection{FILE\_SHARE}
\subsubsection{FILE\_GET}
\subsubsection{FILE\_MODIFY}
\subsubsection{FOLDER\_CAPABILITY}
\subsubsection{FOLDER\_CREATE}
Create a new folder.
What it creates is a function of the task the
server is being asked to do.

Examples of what it could be doing:
\begin{enumerate}
\item When it is processing email messages, it could be creating
a new email folder.
\item When it is processing calendar messages, it could be creating
  a new calendar.
\item Performing actual disk folder (directory) commands.
\end{enumerate}

See the specific task specifications.

A FILE\_CREATE command includes the name of the folder
to create.
And if applicable, a path.
The path separation character is the slash (/) 0x2f UTF-8 character.
\subsubsection{FOLDER\_COPY}
\subsubsection{FOLDER\_DELETE}
\subsubsection{FOLDER\_RENAME}
\subsubsection{FOLDER\_METADATA}
\subsubsection{FOLDER\_MOVE}
\subsubsection{FOLDER\_OPEN}
\subsubsection{FOLDER\_SHARE}
\subsubsection{FOLDER\_LIST}

\subsection{Administrative Commands}
\subsubsection{SERVER\_CONFIGURE}
\subsubsection{SERVER\_KICK\_USER}
\subsubsection{SERVER\_LOGS}
\subsubsection{SERVER\_MANAGE\_BANS}
\subsubsection{SERVER\_SHUTDOWN}
\subsubsection{SERVER\_VIEW\_STATS}
\subsubsection{USER\_CREATE}
\subsubsection{USER\_DELETE}
\subsubsection{USER\_LIST}
\subsubsection{USER\_PERMISSIONS}
\subsubsection{USER\_RENAME}



\chapter{Input / Output Model}
\include{IOModel}
\chapter{Packet Details}
\section{Packet Details}
\subsection{I/O Vector Packet}
\subsection{Comamnd Packet}

\chapter{Code Segments}
\section{Main Thread}\label{sec:MainThread}
The client main thread might be controlled by a user, with
a user interface.
And a client program could exist that automates the
transfer of mime objects without human intervention.

A server is mostly responsive to commands from the client.

Once the main thread has gathered the MIME information to
be transmitted to the server, it takes these basic steps:

\begin{enumerate}
\item Assigns the command (CMD). Commands are defined
  in \ref{todo}.
  
\item Then creates an I/O Vector Packet (\ref{fig:IovPacket}).
  And creates one IO vector Blob for each section of data
  that is part of the command, it appends that Blob object
  to the \_Q member variable in the order needed by CMD.

\item Puts a temporary lock on the output queue \textit{OQ})
  
\item Pushes the new I/O Vector Packet onto the
  end of \textit{OQ}.

\item Removes the temporary lock on the output queue.
  
\end{enumerate}

\section{Dispatch Thread}\label{sec:DispatchThread}


\section{Output Thread}\label{sec:OutputThread}
The output thread ...

\section{XDR Encode}\label{sec:XdrEncode}


\section{TLS Output}\label{sec:TlsIOOut}


\section{TLS Input}\label{sec:TlsIOIn}

\section{XDR Decode}\label{sec:XdrDecode}
It gets packets from the TLS input code (\ref{TlsIOIn}).
As the data is decoded, it is converted from a network
packet to a IovPacket and pushes it into the queue {textit{IQ}).

\begin{figure}
  \centering
  \includesvg{Generated/MultiplePacket.svg}
  \caption{Packet Multiple Commands}
  \label{fig:Packet}
\end{figure}

\chapter{Client Implementation}
\input{ClientImplementation}
\chapter{Server Implementation}
\input{ServerImplementation}
\chapter{Getting the Code, Building, and Testing}
\section{Where to get Phoenix}

\section{Building RPCGEN++}
The entire Phoenix build system is designed to build everything
on Linux for all platforms.
It uses the GNU compilers on Linux.
The mingw compilers to cross compile for Windows.
And the android NDK package to cross compile for Android.
The build has been tested on the Android based Oculus virtual
reality headsets.

The build uses the POSIX make command.
And can build debug, release, and test coverage libraries,
commands, and test code from the make system.

You can build any of the options below, or '\textbf{make all}
to build everything.

For each build, it builds 2 libraries and commands.
First using dynamic shared libraries, or crates them.
And the second pass builds with a minimum number
of dynamic libraries, using as many static libraries
as possible.

The BuildTools directory contains the Makefiles
that are included to do all of the different builds
and cross compiling.

You will need to modify the paths to the mingw and
android NDK paths to match your installation in the
Makefile.top and Makefile.top.lib files.

You will have to set the ANDROID\_CC, ANDROID\_CXX, ANDROID\_STRIP,
W64\_CC, W64\_CXX, and W64\_STRIP settings.

On Linux the mingw packages are often loaded using the operating
systems installation tools (dnf, apt, yum).

The Android NDK packages have to be downloaded using the
latest google instructions, which seem to change over time.
I used the NDK r26 version.

\subsection{Building For Linux}
From the top of the GIT repository you downloaded,
type one of these commands to build for Linux:

\begin{itemize}
\item '\textbf{make LinuxDebug}'
  To build the Linux libraries and commands with symbol
  and debugging information.
  
  The output will be in BuildTemp/Linux/Debug.
  The static library will be in BuildTemp/Linux/Debug/lib
  and the shared library will be named lib{name}.so.{version}.
  The static library will be named lib{name}.a.
  
  Any commands built will be in BuildTemp/Linux/Debug/bin
  for dynamic.

  And BuildTemp/Linux/Debug/sbin for static programs.

\item '\textbf{make LinuxRelease}'
  To build the Linux libraries and commands with the libraries
  and programs compiled with the optimize switch turned on.
  The necessary symbols will be stripped from them to make
  them smaller.

  The output will be in BuildTemp/Linux/Release.
  The static library will be in BuildTemp/Linux/Release/lib
  and the shared library will be named lib{name}.so.{version}.
  The static library will be named lib{name}.a.
  
  Any commands built will be in BuildTemp/Linux/Release/bin
  for dynamic.

  And BuildTemp/Linux/Release/sbin for static programs.

\item '\textbf{make LinuxTcov}'
  To build the Linux libraries and commands with the libraries
  and programs compiled with the test code and test metric
  collection information built into them.
  See test coverage \ref{TestCoverage} for more information.
  Debug symbols are included in TCOV builds.

  The output will be in BuildTemp/Linux/Tcov.
  The static library will be in BuildTemp/Linux/Tcov/lib
  and the shared library will be named lib{name}.so.{version}.
  The static library will be named lib{name}.a.
  
  Any commands built will be in BuildTemp/Linux/Tcov/bin
  for dynamic.

  And BuildTemp/Linux/Tcov/sbin for static programs.

  \textbf{When you run a TCOV build file or library, it generates a LOT of  output files}. Be sure to read \ref{TestCoverage}.
  
\item '\textbf{make Linux}'
  To build LinuxDebug, LinuxRelease, and LinuxTcov.
\end{itemize}

\subsection{Building For Windows on Linux}
  
From the top of the GIT repository you downloaded,
type one of these commands to build for Windows, 64-bit:

\begin{itemize}
\item '\textbf{make W64Debug}'
  To build the Windows 64-bit libraries and commands with symbol
  and debugging information.

  The output will be in BuildTemp/W64/Debug.
  The static library will be in BuildTemp/W64/Debug/lib
  and the shared library will be named lib{name}.so.{version}.
  The static library will be named lib{name}.a.
  
  Any commands built will be in BuildTemp/W65/Debug/bin,
  for dynamic.
  
  And BuildTemp/W64/Debug/sbin for static programs.

\item '\textbf{make W64Release}'
  To build the Windows 64-bit libraries and commands with the libraries
  and programs compiled with the optimize switch turned on.
  The necessary symbols will be stripped from them to make
  them smaller.

  The output will be in BuildTemp/W64/Release.
  The static library will be in BuildTemp/W64/Release/lib
  and the shared library will be named lib{name}.so.{version}.
  The static library will be named lib{name}.a.
  
  Any commands built will be in BuildTemp/W65/Release/bin,
  for dynamic.
  
  And BuildTemp/W64/Release/sbin for static programs.
  
\item '\textbf{make W64Tcov}'
  To build the Windows libraries and commands with the libraries
  and programs compiled with the test code and test metric
  collection information built into them.
  See test coverage \ref{TestCoverage} for more information.
  Debug symbols are included in TCOV builds.

  The output will be in BuildTemp/W64/Tcov.
  The static library will be in BuildTemp/W64/Tcov/lib
  and the shared library will be named lib{name}.so.{version}.
  The static library will be named lib{name}.a.
  
  Any commands built will be in BuildTemp/W64/Tcov/bin
  for dynamic.

  And BuildTemp/W64/Tcov/sbin for static programs.

  \textbf{When you run a TCOV build file or library, it generates a LOT of  output files}. Be sure to read \ref{TestCoverage}.
  
  You have to run the built programs on a Widows device
  with the matching GNU tools.
  
\item '\textbf{make W64}'
  To build W64Debug, W64Release, and W64Tcov.
\end{itemize}

\subsection{Building For Android on Linux}
  
From the top of the GIT repository you downloaded,
type one of these commands to build for Android, 64-bit:
You will need the GCC NDK build tools in your path.

\begin{itemize}
\item '\textbf{make AndroidDebug}'
  To build the Android 64-bit libraries and commands with symbol
  and debugging information.

  The output will be in BuildTemp/Android/Debug.
  The static library will be in BuildTemp/Android/Debug/lib
  and the shared library will be named lib{name}.so.{version}.
  The static library will be named lib{name}.a.
  
  Any commands built will be in BuildTemp/Android/Debug/bin,
  for dynamic.
  
  And BuildTemp/Android/Debug/sbin for static programs.

\item '\textbf{make W64Release}'
  To build the Android 64-bit libraries and commands with the libraries
  and programs compiled with the optimize switch turned on.
  The necessary symbols will be stripped from them to make
  them smaller.

  The output will be in BuildTemp/Android/Release.
  The static library will be in BuildTemp/W64/Release/lib
  and the shared library will be named lib{name}.so.{version}.
  The static library will be named lib{name}.a.
  
  Any commands built will be in BuildTemp/Android/Release/bin,
  for dynamic.
  
  And BuildTemp/Android/Release/sbin for static programs.
  
\item '\textbf{make W64Tcov}'
  To build the Windows libraries and commands with the libraries
  and programs compiled with the test code and test metric
  collection information built into them.
  See test coverage \ref{TestCoverage} for more information.
  Debug symbols are included in TCOV builds.

  The output will be in BuildTemp/Android/Tcov.
  The static library will be in BuildTemp/Android/Tcov/lib
  and the shared library will be named lib{name}.so.{version}.
  The static library will be named lib{name}.a.
  
  Any commands built will be in BuildTemp/Android/Tcov/bin
  for dynamic.

  And BuildTemp/Android/Tcov/sbin for static programs.

  \textbf{When you run a TCOV build file or library, it generates a LOT of  output files}. Be sure to read \ref{TestCoverage}.

  You have to run the built programs on a Android device
  with the matching GNU tools.
  
\item '\textbf{make Android}'
  To build AndroidDebug, AndroidRelease, and AndroidTcov.
\end{itemize}


\subsection{Filing Bugs Reports}
File bug, issues, and questions on github, at:
\href {https://github.com/RiverExplorer/Phoenix/issues}{Bug and Issue Reports}
\subsection{Test Coverage}

\section{Test Coverage}\label{sec:TestCoverage}
This project uses \textit{gcov} and \textit{gcovr} to generate the HTML
test coverage results.

This produces a lot of data and this section will get you
started in understanting what you have produced.

Once the program or library is compiled with the test coverage
tools, you must run them on the target operating system.
So the Windows 64-bit builds (W64) must be run on a windows
computer to get the test results for windows.
And same with the Android devices.

More ....


\chapter{Legal}
\section{Legal Mumbo Jumbo}
\centering{
  \textbf{This document is Copyright 2025 Douglas Mark Royer}.
  }

  \begin{figure}
    \centering
    \includegraphics{RiverExplorerLLC-Logo-BLACK-131x50-Transparent.ps.png}
    \caption{RiverExplorer Company Logo}
    \label{fig:RiverExplorerLogo}
  \end{figure}
  
\centering {
  \textbf{'RiverExplorer' and the logo shown in figure \ref{fig:RiverExplorerLogo} are trademarks of Douglas Mark Royer}.
  \textbf{'RiverExplorer Games LLC' is a wholly owned subsidiary of RiverExplorer LLC.}
  }
\section{License}
\textit{Permission is hereby granted, free of charge, to any person obtaining a copy of this software and associated documentation files (the “Software”), to deal in the Software without restriction, including without limitation the rights to use, copy, modify, merge, publish, distribute, sublicense, and/or sell copies of the Software, and to permit persons to whom the Software is furnished to do so, subject to the following conditions:}

\textit{The above copyright notice and this permission notice shall be included in all copies or substantial portions of the Software.}

\textit{THE SOFTWARE IS PROVIDED “AS IS”, WITHOUT WARRANTY OF ANY KIND, EXPRESS OR IMPLIED, INCLUDING BUT NOT LIMITED TO THE WARRANTIES OF MERCHANTABILITY, FITNESS FOR A PARTICULAR PURPOSE AND NONINFRINGEMENT. IN NO EVENT SHALL THE AUTHORS OR COPYRIGHT HOLDERS BE LIABLE FOR ANY CLAIM, DAMAGES OR OTHER LIABILITY, WHETHER IN AN ACTION OF CONTRACT, TORT OR OTHERWISE, ARISING FROM, OUT OF OR IN CONNECTION WITH THE SOFTWARE OR THE USE OR OTHER DEALINGS IN THE SOFTWARE.}


\appendix
\chapter{Sample XDR files}
\chapter{Sample ABNF files}

\backmatter
\listoffigures
\listoftables
\end{document}
